\documentclass[11pt]{article}

\usepackage[margin=1in]{geometry}
\usepackage{amsmath,amssymb,booktabs}
\usepackage[hidelinks]{hyperref}
\usepackage{microtype}

\title{The Economic Model Behind the Central-Banking Game}
\author{}
\date{}

\newcommand{\gap}{x}
\newcommand{\pistar}{\pi^{*}}
\newcommand{\rstar}{r^{*}}

\begin{document}
\maketitle

\section{Objectives}

The model tries to succeed in two objectives.

First, it should be realistic enough to offer immersion and pedagogical value.
Its equilibria should be reasonable.  In particular, when the real policy rate
is neutral and past policy has also been neutral, output should remain at
potential and inflation should be stable in the long run.  Stable inflation
need not be low inflation: if expected inflation is high but constant, neutral
policy should validate that constant rate rather than cause an artificial
recession.

Second, the model should be recognizable as an undergraduate aggregate-demand
and aggregate-supply model.  Inflation is on the vertical axis and the output
gap is on the horizontal axis.  Short-run aggregate supply slopes upward,
aggregate demand slopes downward, and familiar shocks and policy choices shift
the curves in the expected directions.

These objectives require a distinction between real GDP growth and the output
gap.  Let
\[
  \gap_t \equiv 100\left(\log Y_t-\log Y_t^p\right)
\]
denote the output-level gap.  A positive value means that actual output is
above potential.  If the game reports annualized growth while one turn is one
quarter, actual and potential growth satisfy
\begin{equation}
  g_t=g_t^p+4(\gap_t-\gap_{t-1}).
  \label{eq:growth-gap-identity}
\end{equation}
The AD--AS system determines $(\gap_t,\pi_t)$.  Equation
\eqref{eq:growth-gap-identity} then recovers the GDP growth rate shown to the
player.

\section{Aggregate supply}

Aggregate supply is the simpler curve to motivate.  Following the
expectations-augmented aggregate-supply equation presented by Mishkin and
Serletis in Chapter 23, the model uses
\begin{equation}
  \pi_t=\pi_t^e+\gamma\gap_t+\varepsilon_{\pi,t},
  \qquad \gamma>0.
  \label{eq:as}
\end{equation}
Prices and wages are sticky, so firms respond to unexpectedly strong demand
partly by increasing production and partly by increasing prices.  Equation
\eqref{eq:as} is the output-side counterpart of an expectations-augmented
Phillips curve.

The upward slope is immediate:
\[
  \frac{\partial\pi_t}{\partial\gap_t}=\gamma>0.
\]
A positive supply shock $\varepsilon_{\pi,t}$ shifts AS upward.  Higher
expected inflation also shifts AS upward one for one.

Expected inflation is a weighted average of the central bank's target and
lagged inflation:
\begin{equation}
  \pi_t^e=\alpha_t\pistar_t+(1-\alpha_t)\pi_{t-1},
  \qquad 0\leq\alpha_t\leq1.
  \label{eq:expectations}
\end{equation}
Greater central-bank credibility raises $\alpha_t$ and anchors expectations
more firmly to the target.

\section{Aggregate demand}

\subsection{Why its slope requires care}

Deriving a downward-sloping AD curve is less immediate when the nominal policy
rate is a decision made by the player.  Mishkin and Serletis obtain the usual
negative relationship by endogenizing the central bank's reaction to
inflation.  That approach is not suitable here: the purpose of the game is to
let a central bank decide case by case or create its own policy rule.  Building
a mandatory reaction function into AD would partly make the player's decision
for them.

Instead, the model takes inspiration from Cowen and Tabarrok's use of the
quantity theory of money.  In growth rates, the quantity identity is
\begin{equation}
  m_t+v_t=\pi_t+g_t,
  \label{eq:quantity}
\end{equation}
where $m_t+v_t$ is nominal-demand growth.  Consequently,
\begin{equation}
  g_t=(m_t+v_t)-\pi_t.
  \label{eq:real-growth}
\end{equation}
For a given rate of nominal-demand growth, higher inflation leaves less real
GDP growth.  This provides the negative AD slope without assuming how the
central bank reacts.

\subsection{Nominal-demand growth and monetary policy}

Nominal-demand growth is modeled as
\begin{align}
  m_t+v_t={}&g_t^p+\pi_t^e
  +w_{10}\left[\frac{1}{10}\sum_{j=1}^{10}
    \bigl((i_{t-j}-\pi^e_{t-j})-\rstar_{t-j}\bigr)\right]
  \notag\\
  &+w_{20}\left[\frac{1}{20}\sum_{j=1}^{20}
    \bigl((i_{t-j}-\pi^e_{t-j})-\rstar_{t-j}\bigr)\right]
  \notag\\
  &+\beta_r(i_t-\pi_t^e-\rstar_t)+\varepsilon_{d,t},
  \label{eq:nominal-demand}
\end{align}
with $w_{10}<0$, $w_{20}<0$, and $\beta_r<0$.  When fewer observations are
available, each average uses all available history.

Every real-rate gap is ex ante: the nominal rate is paired with the inflation
expectation prevailing when that rate was chosen.  The game therefore records
past expected inflation as part of its quarterly history rather than rebuilding
past real rates from realized inflation.  A real rate above $\rstar$ reduces
nominal-demand growth; a rate below $\rstar$ raises it.  The moving averages
represent monetary-policy lags.

The term $\pi_t^e$ deserves emphasis.  It is not added merely to force a chosen
inflation target into AD.  Neutral nominal spending must grow quickly enough to
support both potential real growth and the inflation rate already incorporated
in nominal plans and contracts.  Thus neutral nominal-demand growth is
\[
  g_t^p+\pi_t^e.
\]
Without $\pi_t^e$, positive steady inflation would mechanically subtract from
real growth in equation \eqref{eq:real-growth}.  Even with neutral interest
rates, the model would push output below potential.  Including $\pi_t^e$ is
therefore necessary for the desired equilibrium in which $\pi_t=\pi_t^e$ and
$\gap_t=0$.

\subsection{AD in output-gap space}

Equating \eqref{eq:growth-gap-identity} and \eqref{eq:real-growth} gives
\[
  g_t^p+4(\gap_t-\gap_{t-1})=(m_t+v_t)-\pi_t.
\]
Substituting the complete nominal-demand equation \eqref{eq:nominal-demand}
and cancelling $g_t^p$ yields the AD curve:
\begin{align}
  \gap_t={}&\gap_{t-1}+\frac{1}{4}\Bigg\{
  \pi_t^e-\pi_t
  +w_{10}\left[\frac{1}{10}\sum_{j=1}^{10}
    \bigl((i_{t-j}-\pi^e_{t-j})-\rstar_{t-j}\bigr)\right]
  \notag\\
  &+w_{20}\left[\frac{1}{20}\sum_{j=1}^{20}
    \bigl((i_{t-j}-\pi^e_{t-j})-\rstar_{t-j}\bigr)\right]
  +\beta_r(i_t-\pi_t^e-\rstar_t)+\varepsilon_{d,t}
  \Bigg\}.
  \label{eq:ad}
\end{align}

All interest-rate variables in equation \eqref{eq:ad} shift AD.  Holding them
and the inherited gap fixed,
\[
  \frac{\partial\gap_t}{\partial\pi_t}=-\frac{1}{4}<0.
\]
AD therefore slopes downward without embedding a central-bank reaction rule.
Higher current nominal rates shift AD left because
\[
  \frac{\partial\gap_t}{\partial i_t}=\frac{\beta_r}{4}<0.
\]
Restrictive historical real-rate gaps also shift AD left because both history
weights are negative.

\section{Equilibrium}

On the ordinary upward-sloping portion of AS, equations \eqref{eq:as} and
\eqref{eq:ad} jointly determine current inflation and the output gap.  Consider
the neutral case:
\begin{itemize}
  \item $\gap_{t-1}=0$;
  \item all historical real-rate gaps are zero;
  \item $i_t-\pi_t^e=\rstar_t$;
  \item both shocks are zero.
\end{itemize}
Equation \eqref{eq:ad} reduces to
\[
  \gap_t=\frac{1}{4}(\pi_t^e-\pi_t),
\]
while AS is
\[
  \pi_t=\pi_t^e+\gamma\gap_t.
\]
Their unique intersection is
\[
  \boxed{\gap_t=0,\qquad \pi_t=\pi_t^e.}
\]
Equation \eqref{eq:growth-gap-identity} then gives $g_t=g_t^p$.  The equilibrium
works for any stable expected inflation rate, not only for the target rate.

\section{Unemployment and capacity}

Okun's law translates the output-level gap into an unemployment gap:
\begin{equation}
  u_t=u_t^n-\beta_u\gap_t,
  \qquad \beta_u>0.
  \label{eq:okun}
\end{equation}
Because $\gap_t$ is a level gap, unemployment inherits the accumulated effects
of past above- or below-potential growth.  It no longer resets each quarter as
a function of current growth alone.  At $\gap_t=0$, unemployment equals its
natural rate.

The upward-sloping AS segment ends at a capacity gap inferred from the minimum
sustainable unemployment rate $u^{\min}$:
\begin{equation}
  \gap^{\max}=\frac{u_t^n-u^{\min}}{\beta_u}.
  \label{eq:capacity}
\end{equation}
Beyond this point AS is vertical.  If unconstrained AD calls for
$\gap_t>\gap^{\max}$, output is held at $\gap^{\max}$ and the downward-sloping
AD curve determines inflation.  Moving left along a downward-sloping AD curve
raises inflation, so the constrained solution lies above the AS kink as a
textbook diagram requires.

\section{Variables reported to the player}

The equilibrium variables are $\gap_t$ and $\pi_t$.  After they are solved, the
game reports
\begin{align*}
  g_t&=g_t^p+4(\gap_t-\gap_{t-1}),\\
  u_t&=u_t^n-\beta_u\gap_t.
\end{align*}
Potential growth is therefore needed to report actual GDP growth, but it drops
out of the cyclical AD--AS equilibrium.  This separation keeps the diagram
textbook-like while retaining a familiar growth statistic in the interface.

\begin{thebibliography}{9}
\bibitem{mishkin-serletis}
Frederic S. Mishkin and Apostolos Serletis,
\emph{The Economics of Money, Banking, and Financial Markets}, Chapter 23.

\bibitem{cowen-tabarrok}
Tyler Cowen and Alex Tabarrok,
\emph{Modern Principles of Economics}, discussion of aggregate demand and the
quantity theory of money.
\end{thebibliography}

\end{document}
